\chapter{Repeated decisions and customers who disappear}
\label{gap29:servicing}

A servicing policy changes what happens next. A reminder, payment holiday or
collection action affects both the customer's state and the information the
bank subsequently observes. Estimating a policy from customers who remain in
the records can therefore answer a different question from the one the bank
faces. The relevant value averages over the initial eligible population,
including customers whose later outcomes are missing.

An explicitly generated study makes this distinction measurable. With a
correct observation model, weighting recovers the chosen policy's value on
average, but repeated actions and dropout leave little effective information.
When dropout depends on an unobserved future shock, the same calculation
becomes biased and its nominal 95\% interval covers the numerical reference
in only 80\% of replications. Analysing complete observed trajectories also
overstates value. These are findings under the stated constructed law,
not estimates of a bank's servicing policy.

\section{A policy changes a path, not just a row}
Each generated customer has a fixed characteristic $X\sim N(0,1)$ and initial
state $L_0=0.4X+\varepsilon_0$. There are three decisions, $t=0,1,2$.
The historical action is binary, with probability
\[
 p_t=\operatorname{clip}\{\operatorname{expit}(0.3X-0.5L_t+0.2A_{t-1}),0.2,0.8\},
 \qquad A_{-1}=0.
\]
Clipping is a feature of this generated assignment law, not a silent
alteration of weights estimated from observed bank data. The next state and
reward are
\[
 L_{t+1}=0.6L_t+0.25X-0.5A_t+\varepsilon_{t+1},\qquad
 R_t=10-2L_{t+1}^2-A_t+0.5A_tX,
\]
with independent standard-normal innovations. Actions have a direct cost
and change the future state. Rewards are illustrative units; the coefficients
have not been estimated from a bank.

Observation continues with probability
$g_t=0.65+0.30\operatorname{expit}(0.3L_t-0.2A_t)$ after the action but
before the next state and reward are observed. Dropout is absorbing in the
analyst's data. The baseline observation draw is independent of the next
innovation conditional on the known history and action. The adverse case
instead uses $\operatorname{expit}\{\operatorname{logit}(g_t)+1.5\varepsilon_{t+1}\}$.
The analyst cannot observe that innovation when specifying the observation
probability. This deliberately violates the assumption behind the weighting
and regression procedures.

The policy value is $V(\pi)=E_\pi[\sum_{t=0}^2R_t]$. The comparison includes
never acting and four threshold policies, which act when $L_t$ exceeds
$-0.5$, $0$, $0.5$ or $1$. A selected policy chooses among the thresholds
using 2,000 separate training trajectories; its evaluation uses 1,000 new
trajectories. Selection and evaluation therefore do not reuse outcomes.
The finite list is a meaningful alternative set, not the class of all
possible policies.

\section{The denominator must retain the initial customers}
For each observed reward, the prefix-weighted score is
\[
 S_i(\pi)=\sum_{t=0}^2 R_{it}
  \prod_{s=0}^t
  \frac{\mathbf1\{A_{is}=\pi(H_{is})\}\mathbf1\{C_{is}=1\}}
       {p_s(A_{is}\mid H_{is})g_s(H_{is},A_{is})}.
\]
The estimate averages $S_i$ over all initial customers. Successive conditional
expectations remove the observation and action probabilities when those
probabilities are correct, positive and based on sufficient observed history.
The reward-specific prefix is necessary: observing a reward at the last date
requires surviving every preceding observation decision.

The simple competing calculation averages total reward among customers whose
whole recorded action sequence matches the policy and whose records remain
complete. Its denominator selects customers. It has no reason to equal
$V(\pi)$. A fitted alternative estimates the state transition from observed
transitions and simulates policy paths using that regression. This
g-computation exercise receives the constructed reward rule and initial-state
law. Its smaller observed errors in the correctly specified examples do not
make it a ready-made bank estimator. Parameter uncertainty is not propagated
into an interval for this fitted alternative in the first comparison.
Chapter~\ref{gap29:policy-uncertainty} then examines that missing uncertainty
using whole-customer resampling and the same adverse mechanisms.

The study runs 200 independent replications in each of four cases: the
baseline; action probabilities allowed to range from 0.02 to 0.98; dropout
depending on the future innovation; and omission of $X$ from the fitted state
regression. Each policy also has a separate 100,000-trajectory numerical
reference, with its own Monte Carlo standard error. Reference trajectories
are unavailable to selection and estimation.

\section{A usable interval requires enough information}
For the fixed threshold 0.5 policy, the reference value is 20.250 with Monte
Carlo standard error 0.030. Never acting has value 19.493; their paired
reference difference is 0.756 with Monte Carlo standard error 0.014. In the
baseline, weighting has estimated bias -0.035 with Monte Carlo standard
error 0.142. The complete-case calculation overstates value by 1.915
(Monte Carlo standard error 0.083). This distortion comes from selection,
not from a language or interface difference.

\begin{table}[htbp]
\centering\small
\begin{tabular}{lrrrr}
\toprule
Case & Weighted bias & Complete-case bias & Fitted bias & Coverage \\
\midrule
Baseline & -0.035 & 1.915 & -0.006 & 0.950 \\
Wider assignment range & -0.053 & 1.987 & -0.025 & 0.950 \\
Future-dependent dropout & -1.874 & 3.036 & 1.773 & 0.800 \\
Omitted state predictor & 0.117 & 1.872 & -0.281 & 0.930 \\
\bottomrule
\end{tabular}
\caption{Generated threshold 0.5 policy, 200 replications per case. Coverage
is for the normal prefix-weighted interval against the numerical reference;
fitted g-computation has no reported confidence interval. Bias is in
illustrative reward units, not currency.}
\end{table}

Baseline weighted coverage is 0.950, with a Monte Carlo standard error of
0.015 and a Wilson interval approximately [0.910, 0.973]. This is compatible
with the nominal rate, not a precise calibration certificate. Average
effective sample sizes fall from about 320 at the first reward to 97 and
29 at the later rewards. The average interval width is 8.047, far larger
than the policy improvement of 0.756. A positive lower bound for the paired
improvement appears in only 3.5\% of baseline replications. Correct weighting
therefore leaves an information problem even in this favourable construction.

With future-dependent dropout, the weighted bias is -1.874 (Monte Carlo
standard error 0.125), and coverage falls to 0.800. Its Monte Carlo Wilson
interval is approximately [0.739, 0.850]. Allowing for reference values within
their separate 95\% numerical band gives coverage bounds of [0.790, 0.820].
The failure is not explained by reference simulation noise. Fitted
g-computation is biased in the opposite direction, by about 1.773. The
observation problem affects both approaches.

When the state regression omits $X$, its estimated bias is -0.281 (Monte
Carlo standard error 0.033). Weighting still uses the correct observed-history
probabilities in that case. This is a useful distinction: a model can fail
because its transition equation is wrong, because its observation mechanism
is wrong, or because too few informative trajectories remain. These failures
call for different repairs.

The separately selected threshold has average regret about 0.248 relative
to the finite candidate list in the baseline. Selection is therefore not
equivalent to finding the best policy even when evaluation is kept separate.
The reference values nominate threshold 1 within this list; they establish
no universal policy optimum. This study does not rank viable estimators by
a few realised errors or declare that a single fitted model is generally
superior.

\section{What the next prototype must learn}
The bank needs to know why a customer's record ends, what was available at
each action, and whether all relevant actions had a reasonable chance of
being taken. More complete outcome follow-up or an audit sample can address
missing information. More representative action variation can improve
support. A better state model can help prediction, but it cannot by itself
recover outcomes selected by an unobserved future shock.

The software implication is concrete. Preserve the initial denominator,
reward-specific weights, information dates and policy-selection split; show
effective information and sensitivity alongside the effect. A commercial or
open-source implementation should be compared on these same meanings. A
successfully executed policy-value calculation does not settle the scientific
conditions that make it useful. Here the next experiment should test those
conditions before seeking a costly bank pilot.
